% Generated draft: requires mathematical and source review.
\documentclass[11pt]{article}
\usepackage[margin=25mm]{geometry}
\usepackage{fontspec}
\setmainfont{Latin Modern Roman}
\usepackage{amsmath,amssymb,mathtools,unicode-math}
\setmathfont{Latin Modern Math}
\usepackage{enumitem,needspace,xurl,hyperref}
\setlength{\parindent}{0pt}
\setlength{\parskip}{5pt}
\setlength{\emergencystretch}{3em}
\newcommand{\N}{\mathbb{N}}
\newcommand{\Z}{\mathbb{Z}}
\newcommand{\Q}{\mathbb{Q}}
\newcommand{\R}{\mathbb{R}}
\newcommand{\C}{\mathbb{C}}
\newcommand{\F}{\mathbb{F}}
\newcommand{\A}{\mathbb{A}}
\newcommand{\PP}{\mathbb{P}}
\newcommand{\E}{\mathbb{E}}
\newcommand{\eps}{\varepsilon}
\newcommand{\dd}{\,\mathrm{d}}
\begin{document}
\section*{UnsolvedMath ID: 30004331}
\textbf{Source problem code:} catalogue entry
\subsection{Definitions and mathematical statement}
Interpret a complex projective line arrangement as a reduced union
\[
C=L_1\cup\cdots\cup L_d\subset\PP^2_{\C}
\]
of finitely many distinct lines. Let \(S=\operatorname{Sing}(C)\), let \(s=|S|\), and let \(m_P\) be the number of lines through \(P\in S\). For \(s>0\), let \(\pi:X\to\PP^2_{\C}\) be the blow-up at all points of \(S\), and let \(\widetilde C\) be the strict transform of \(C\). Define
\[
H(C)=\frac{\widetilde C^{\,2}}{s}
     =\frac{d^2-\sum_{P\in S}m_P^2}{s}.
\]
The first question is whether every such complex line arrangement satisfies
\[
H(C)\ge -\frac{225}{67}.
\]
The second question asks for
\[
\inf\{H(C):C\text{ is a rational line arrangement},\ |\operatorname{Sing}(C)|>0\}.
\]
The source does not define ``rational line arrangement.'' If it means an arrangement whose individual lines are defined over \(\Q\), this asks for the infimum over those arrangements, viewed after extension of scalars to \(\C\). This interpretation requires confirmation: rationality of the components alone imposes no restriction on projective line arrangements. The exclusion \(s=0\) is necessary for the displayed definition of \(H(C)\).

For the ideal-theoretic questions, fix a ground field \(k\) and a nonempty finite reduced set of closed points \(Z\subset\PP^2_k\). Write \(R=k[x_0,x_1,x_2]\), let \(\mathfrak p_P\subset R\) be the homogeneous prime ideal of \(P\), and define
\[
I_Z=\bigcap_{P\in Z}\mathfrak p_P,
\qquad
I_Z^{(m)}=\bigcap_{P\in Z}\bigl(\mathfrak p_P^mR_{\mathfrak p_P}\cap R\bigr)
\quad(m\ge1).
\]
Here \(I_Z^r\) denotes the ordinary ideal power. The third question is to characterize all finite point sets in the intended ground-field setting for which
\[
I_Z^{(3)}\not\subseteq I_Z^2.
\]
The source does not specify the ground field, its characteristic, or whether points must be rational over that field; the definition above gives a field-relative formulation without resolving that missing scope. A complex-ground-field interpretation is compatible with the preceding arrangement question but is not explicitly stated for this question. If the empty set is allowed, its ideal is \(R\), so it does not satisfy the displayed noncontainment.

For a nonzero proper point ideal \(I\subset R\), define its resurgence by
\[
\rho(I)=\sup\left\{\frac{m}{r}:m,r\in\Z_{\ge1},\ I^{(m)}\not\subseteq I^r\right\}.
\]
For a line arrangement \(C\) with nonempty singular set, put \(I=I_{\operatorname{Sing}(C)}\). The fourth question asks to determine \(\rho(I)\) for line arrangements, at least for supersolvable arrangements. Here supersolvable has its usual arrangement-theoretic meaning: the intersection lattice of the associated central arrangement in \(k^3\) has a maximal chain of modular elements. In its ranked lattice, an element \(X\) is modular when
\[
\operatorname{rk}(X)+\operatorname{rk}(Y)
=\operatorname{rk}(X\wedge Y)+\operatorname{rk}(X\vee Y)
\quad\text{for every }Y.
\]
The source supplies no particular formula, algorithmic specification, or additional restriction for this requested computation. The intended ground-field scope of this last question also requires confirmation.
\subsection{Short English statement}
The record asks whether every complex line arrangement with a defined Harbourne constant has \(H(C)\ge-225/67\), and asks for the infimum over ``rational'' line arrangements. It also asks which finite plane point sets violate \(I_Z^{(3)}\subseteq I_Z^2\), and seeks the resurgence of singular-point ideals of line arrangements, at least in the supersolvable case. The meaning of ``rational'' and the ground fields for the ideal questions remain unspecified. The supplied open status is dated curation, not an independently verified conclusion.
\subsection{Sources}
\begin{itemize}
\item Local UnsolvedMath registry: permanent ID \(30004331\), source problem code \texttt{catalogue entry}, title ``Harbourne Constants, Symbolic Powers, and Resurgence of Line Arrangements,'' in the supplied Oberwolfach Reports Open Problems record.
\item Supplied source citation: \emph{Mini-Workshop: Seshadri Constants (2020)}, Oberwolfach Reports, pp.\ 3271--3313, DOI \texttt{10.4171/owr/2019/53}. \url{https://doi.org/10.4171/owr/2019/53}. The registry records proposed year 2019 and names no individual proposer. The relationship between that year, the DOI year, and the supplied 2020 citation has not been independently checked.
\item Supplied literature source: ``OWR report on Seshadri constants.'' \url{https://ems.press/content/serial-article-files/46833}.
\item Supplied literature source: ``Bauer et al., bounded negativity and arrangements of lines.'' \url{https://arxiv.org/abs/1407.2966}.
\item Literature assessment dated exactly \texttt{2026-08-22}: the supplied curation labels the record open and states that subsequent arrangement results give examples and bounds but do not determine the requested infima or characterize all containments. This is supplied dated literature triage derived from the OWR curation, not an independently verified assessment; its time-sensitive status and cited claims require independent verification.
\item The registry labels the statement \texttt{corrected\_verified}, while its verification note expressly limits the repairs to formula delimiters, notation, grammar, and typography, preserving mathematical content and scope. This label does not establish mathematical review of this draft.
\end{itemize}
\end{document}
